% Generated draft: requires mathematical and source review.
\documentclass[11pt]{article}
\usepackage[margin=25mm]{geometry}
\usepackage{fontspec}
\setmainfont{Latin Modern Roman}
\usepackage{amsmath,amssymb,mathtools,unicode-math}
\setmathfont{Latin Modern Math}
\usepackage{enumitem,needspace,xurl,hyperref}
\setlength{\parindent}{0pt}
\setlength{\parskip}{5pt}
\setlength{\emergencystretch}{3em}
\newcommand{\N}{\mathbb{N}}
\newcommand{\Z}{\mathbb{Z}}
\newcommand{\Q}{\mathbb{Q}}
\newcommand{\R}{\mathbb{R}}
\newcommand{\C}{\mathbb{C}}
\newcommand{\F}{\mathbb{F}}
\newcommand{\A}{\mathbb{A}}
\newcommand{\PP}{\mathbb{P}}
\newcommand{\E}{\mathbb{E}}
\newcommand{\eps}{\varepsilon}
\newcommand{\dd}{\,\mathrm{d}}
\begin{document}
\section*{UnsolvedMath ID: 30001698}
\textbf{Source problem code:} catalogue entry
\subsection{Definitions and mathematical statement}
A finite abstract simplicial complex on a finite vertex set \(V\) is a family \(K\subseteq 2^V\) containing \(\varnothing\) and every singleton \(\{v\}\), and satisfying \(\tau\subseteq\sigma\in K\Rightarrow\tau\in K\). Its members are called faces. The complex is \emph{flag} if every subset \(S\subseteq V\) whose two-element subsets are all faces of \(K\) is itself a face of \(K\).

For \(V\ne\varnothing\), define
\[
\dim K=\max_{\sigma\in K}(|\sigma|-1),\qquad
f_i(K)=\bigl|\{\sigma\in K:|\sigma|=i+1\}\bigr|\quad(0\le i\le\dim K).
\]
The face vector under the convention used here is \(f(K)=(f_0(K),\ldots,f_{\dim K}(K))\); the empty face is not a coordinate.

\textbf{Characterization problem.} For every integer \(d\ge0\), give explicit necessary and sufficient conditions on an integer vector \(a=(a_0,\ldots,a_d)\in\Z^{d+1}\) for the existence of a finite flag simplicial complex \(K\) such that
\[
\dim K=d\quad\text{and}\quad f_i(K)=a_i\quad\text{for every }0\le i\le d.
\]

\textbf{Unresolved source qualifications.} The supplied statement does not specify face-vector indexing, the treatment of complexes without vertices, or trailing zero coordinates. The formulation above makes explicit the convention of recording all nonempty-face counts through the actual dimension for complexes with at least one vertex. If the empty-face coordinate is included, its value under these definitions is \(f_{-1}=1\). Recording a finite integer value for \(f_0\) already forces the vertex set, and hence the complex, to be finite. The supplied background describes an incomplete construction discussion, but supplies neither that construction nor additional hypotheses; no such hypotheses are inferred here.
\subsection{Short English statement}
Characterize exactly the integer vectors that count faces of each dimension in finite flag simplicial complexes. Here flag means that every set of pairwise adjacent vertices forms a face. The source leaves indexing and degenerate cases unspecified; the explicit convention above requires review.
\subsection{Sources}
\begin{itemize}
\item Local registry: \texttt{data/unsolvedmath/problems.json}, permanent UnsolvedMath ID \(30001698\), source problem code \texttt{catalogue entry}, title \emph{Face Vectors of Flag Simplicial Complexes}; collection: Oberwolfach Reports Open Problems. The supplied proposed year is 2011; no individual proposer is identified.
\item Supplied report citation and literature-source link: \emph{Topological and Geometric Combinatorics} (2011), Oberwolfach Reports, pp. 349--423, DOI \texttt{10.4171/owr/2011/08}, \url{https://doi.org/10.4171/owr/2011/08}. The supplied background describes these reports as documenting workshop discussions and open problems. No pinpoint location within the report is supplied.
\item Supplied literature assessment, checked on \textbf{2026-08-21}: the record labels the problem open and describes an incomplete discussion of a construction for flag-complex face vectors that states no resolution of the characterization problem. This is supplied dated curation derived from the OWR material, not an independently verified literature assessment. The supplied source checked is the original Oberwolfach report linked above; time-sensitive status and cited claims require independent verification.
\item The registry labels the statement exact and reports that manual review found the extracted statement self-contained, mathematically coherent, and free of extraction artifacts. These are supplied verification claims; the report has not been independently consulted for this draft.
\end{itemize}
\end{document}
